% Teil II: der durchgerechnete Weg.

% Gedankenschritt: Quadratische Gleichungen über Faktorisierung oder
% Formel; Wurzelgleichungen quadrieren und zwingend zurücksetzen (Probe).

\subsection*{Gegeben und gesucht}
Quadratische Gleichungen (a)--(c) und Wurzelgleichungen (d)--(g) nach der
reellen Zahl $x$ auflösen.

\subsection*{Ansatz}
Quadratische: auf die Form $=0$ bringen, dann faktorisieren oder Formel.
Wurzelgleichungen: Wurzel isolieren, quadrieren, lösen und
\emph{immer} die ursprüngliche Gleichung als Probe rechnen.

\subsection*{a) $x^2+4x=5$}
\begin{equation}
  x^2+4x-5=0 \Rightarrow (x+5)(x-1)=0 \Rightarrow x=-5 \text{ oder } x=1.
\end{equation}
Nullstellen-Suche: Membran der Faktorisierung ist $(x+5)(x-1)=x^2-x+5x-5$.

\subsection*{b) $2x^2+8x-10=32$}
\begin{equation}
  2x^2+8x-42=0 \Rightarrow x^2+4x-21=0\Rightarrow (x+7)(x-3)=0
  \Rightarrow x=-7 \text{ oder } x=3.
\end{equation}
(Nach Division durch $2$ gilt die Faktorisierung weiter.)

\subsection*{c) $x^2+3x=4$}
\begin{equation}
  x^2+3x-4=0\Rightarrow (x+4)(x-1)=0\Rightarrow x=-4 \text{ oder } x=1.
\end{equation}

\subsection*{d) $\sqrt{x+2}=3$}
Hier ist das Quadrieren \emph{unbedingt} legitim, weil beide Seiten der
Gleichung offensichtlich $\ge0$ sind:
\begin{equation}
  x+2=9\Rightarrow x=7.
\end{equation}
Probe: $\sqrt{7+2}=\sqrt9=3$. \checkmark

\subsection*{e) $\sqrt{x+1}=x-5$}
Quadratieren: $x+1=(x-5)^2=x^2-10x+25$, also $x^2-11x+24=0$, d.h.
$x=3$ oder $x=8$. \textbf{Probe:} $x=3$: $\sqrt4=2$, rechts $-2$ --
eine Wurzel ist niemals negativ, deshalb \emph{verwerfen};  $x=8$:
$\sqrt9=3$, rechts $3$ -- \emph{gültig}.
\begin{equation}
  x=8 .
\end{equation}

\subsection*{f) $\sqrt{x+4}=\sqrt{2x+3}$}
Quadratieren: $x+4=2x+3\Rightarrow x=1$. Probe: $\sqrt5=\sqrt5$. \emph{gültig}.

\subsection*{g) $\sqrt{x+1}+\sqrt{x+5}=4$}
\textbf{Gedankenschritt:} Hier reicht ein einfaches Quadrieren nicht, es
bleiben Wurzeln übrig. Statt stumpf weiterzuquadrieren, setzen wir
$a=\sqrt{x+1}$, $b=\sqrt{x+5}$. Dann ist $a+b=4$ und
$b^2-a^2=(x+5)-(x+1)=4$; die dritte binomische Formel gibt
$(b-a)(b+a)=4\Rightarrow b-a=1$. Zusammen $a+b=4$ und $b-a=1$:
$a=\tfrac32$, $b=\tfrac52$, damit $x+1=\tfrac94$, also
\begin{equation}
  x=\frac54 .
\end{equation}
Probe: $\sqrt{\tfrac54+1}=\sqrt{\tfrac94}=\tfrac32$,
$\sqrt{\tfrac54+5}=\sqrt{\tfrac{25}{4}}=\tfrac52$, Summe $4$. \emph{gültig}.

\begin{ergebnis}
\begin{align*}
  \text{a)}\ x\in\{-5,1\}, &\quad
  \text{b)}\ x\in\{-7,3\}, &\quad
  \text{c)}\ x\in\{-4,1\}, &\quad
  \text{d)}\ x=7,\\
  \text{e)}\ x=8, &\quad
  \text{f)}\ x=1, &\quad
  \text{g)}\ x=\tfrac54 .
\end{align*}
\end{ergebnis}

\begin{kontrolle}
\textbf{Probe in der Originalgleichung (zwingend):} Die Quadrierung
lieferte bei e) zusätzlich $x=3$; die Rückprobe zeigt $2\neq-2$ und
verwirft sie --- \emph{nur} so bleibt $x=8$ übrig. In d), f), g) ist die
Probe ebenfalls erfolgreich.

\textbf{Grenzfall / Plausibilität:} g) Annäherung $x=\tfrac54=1{,}25$ ist
im Definitionsbereich ($x\ge-1$ und $x\ge-5$), die Wurzeln sind
$\tfrac32$ und $\tfrac52$, Summe $4$, weiter nichts. Wäre das errechnete
$a$ negativ gewesen, hätte man es verworfen, weil die Wurzel nicht negativ
sein kann.
\end{kontrolle}

\begin{merksatz}
Quadratische Gleichungen per Faktorisierung oder Formel; Wurzelgleichungen
immer quadrieren und \emph{proben}, weil das Quadrieren Scheinlösungen
erzeugen kann. Bei zwei Wurzeln: Substitution $a,b$ + Differenz der
Wurzelquadrate (dritte binomische Formel) ist oft der pfiffige Weg.
\end{merksatz}
